%% kapa: exam=1 round=28 year=2017 period=2 form=A subject=accounting subject_ko=회계학 seq=20 number=60
%% source: Q-Net 감정평가사 기출문제 (cst00309) artlSeq=5199603 file=extracted/1008/5199603/2017년 제28회 감정평가사 1차시험 문제지/2교시 A형.pdf
%% conversion: kordoc 4.24.51 -> tools/gen1.py (auto); review=proofread-sonnet; flags=none
\begin{question}{60}
회계정책, 회계추정의 변경 및 오류에 관한 설명으로 옳은 것은?
\begin{choices}
  \item 측정기준의 변경은 회계정책의 변경이 아니라 회계추정의 변경에 해당한다.
  \item 회계추정의 변경효과를 전진적으로 인식하는 것은 추정의 변경을 그것이 발생한 시점 이후부터 거래, 기타 사건 및 상황에 적용하는 것을 말한다.
  \item 과거에 발생한 거래와 실질이 다른 거래, 기타 사건 또는 상황에 대하여 다른 회계정책을 적용하는 경우에도 회계정책의 변경에 해당한다.
  \item 과거기간의 금액을 수정하는 경우 과거기간에 인식, 측정, 공시된 금액을 추정함에 있어 사후에 인지된 사실을 이용할 수 있다.
  \item 회계정책의 변경과 회계추정의 변경을 구분하는 것이 어려운 경우에는 이를 회계정책의 변경으로 본다.
\end{choices}
\end{question}
\answer{2}
